\documentclass[11pt]{article}
\usepackage[margin=1in]{geometry}
\usepackage{graphicx}
\usepackage{booktabs}
\usepackage{array}
\usepackage{amssymb}
\usepackage{hyperref}
\usepackage{xcolor}
\usepackage{longtable}
\usepackage{caption}
\hypersetup{colorlinks=true,linkcolor=blue,citecolor=blue,urlcolor=blue}
\newcommand{\SUtwoD}{\ensuremath{SU(2)_D}}
\newcommand{\gev}{\,{\rm GeV}}
\newcommand{\msun}{M_\odot}
\title{Final Report: Minimal Conformal Non-Abelian Model for a PTA-Scale Signal}
\author{DarkAgents}
\date{2026-06-07}

\begin{document}
\maketitle

\section{Summary}
The completed \texttt{fopt-pta} pipeline identifies the minimal conformal non-Abelian candidate as a classically scale-invariant dark \SUtwoD{} gauge theory with one complex scalar doublet. The local authoritative model uses a zero-portal secluded dark sector, two free parameters $(g_D,v_D)$, and the vector-mass convention $m_{W_D}(\chi)=g_D\chi/2$.

The exact core model is not new: Borah, Dasgupta and Kang studied the same classically conformal dark \SUtwoD{} doublet model as a NANOGrav 12.5-year first-order phase-transition interpretation in Ref.~\cite{Borah2021}. The present pipeline should therefore be read as a local recomputation and updated PTA-likelihood exercise, not as a model-level novelty claim.

The first-order phase-transition (FOPT) scan found many viable points. The PTA benchmark scan found a near-match to the PTA violin windows, with best score $13/14$. The PTArcade campaign completed with status \texttt{ok}, $N_{\rm samples}=100000$, and maximum-a-posteriori (MAP) point
\[
g_D=0.863175,\qquad v_D=1.10802\gev .
\]
This is a conditional result. The dominant limitations are the DBF gravitational-wave template validity at extremely large $\alpha$, the absence of an independent wall-velocity or efficiency calculation, and missing thermal-history, BBN/CMB, $\Delta N_{\rm eff}$, relic-density, and self-interaction observables. The constraints stage found no direct exclusion, but mostly because the necessary observables are absent, not because full cosmological viability has been proven.

\section{Prompt}
\begin{quote}
what is the minimal conformal non-abelian model that can explain the NANOGrav                     
  signal?
\end{quote}

\section{Model Description}
The critic-authoritative model is \texttt{minimal\_conformal\_su2\_doublet}. It contains a dark non-Abelian gauge group \SUtwoD{} with coupling $g_D$ and one complex scalar doublet $\Phi$, while all new fields are Standard-Model singlets. Classical scale invariance is imposed at tree level, forbidding explicit scalar mass terms. In unitary gauge,
\[
\Phi=\begin{pmatrix}0\\ \chi/\sqrt{2}\end{pmatrix},\qquad \langle \chi\rangle=v_D .
\]
The tree-level scalar potential is
\[
V_0(\chi)=\frac{\lambda_\phi}{4}\chi^4,
\]
where $\lambda_\phi$ is Coleman-Weinberg-fixed rather than an independent scan parameter. The symmetry-breaking pattern is
\[
\SUtwoD \rightarrow \varnothing,
\]
so the three dark vectors become massive, with
\[
m^2_{W_D}(\chi)=\frac{g_D^2\chi^2}{4}.
\]
The free parameters used downstream are $g_D>0$ and $v_D>0$. The local handoff flags a convention tension with Ref.~\cite{Borah2021}: that paper states a physical vector mass convention consistent with $M_{Z_D}=g_D M/2$, while also displaying a field-dependent mass normalization that is not directly identical to the local convention. All local computations use the critic-authoritative convention above.

\section{Pipeline Status}
\begin{center}
\begin{tabular}{p{0.22\linewidth}p{0.14\linewidth}p{0.52\linewidth}}
\toprule
Stage & Status & Interpretation \\
\midrule
Model / critique & warning & Usable model; no red flags, but same-model precedent and mass-convention caveat. \\
Preliminary literature & warning & Exact core model already studied in Ref.~\cite{Borah2021}; benchmark conventions require care. \\
FOPT scan & ok & 934 points scanned, 662 viable, 268 failed, 4 excluded. \\
PTA analysis & warning & Benchmark and PTArcade completed, but no $14/14$ benchmark and template/prior caveats remain. \\
Constraints & warning & No exclusion claimed; most relevant observables require new backend calculations. \\
Prior audit & warning & High-severity limitations dominate the final interpretation. \\
\bottomrule
\end{tabular}
\end{center}

\section{Preliminary Literature Search}
The same-model precedent is Ref.~\cite{Borah2021}, which studied a classically conformal dark \SUtwoD{} model with one complex scalar doublet and a negligible Higgs portal for a NANOGrav 12.5-year FOPT interpretation. The directly quoted benchmark from that work is $g_D=1.37$, $M_{Z_D}=8.23\,{\rm MeV}$, $T_\star=1.8\,{\rm MeV}$, $\alpha_\star=0.45$, and $\phi_c/T_c=3.56$. This benchmark should not be imported into the local parameter basis without resolving the vector-mass convention.

Several structurally related works provide context rather than direct benchmarks. Classically conformal Abelian hidden-sector models were studied for nano-Hertz gravitational waves in Refs.~\cite{Balan2025,Goncalves2025}. Nearly conformal or phenomenological dark-sector phase-transition explanations of NANOGrav were considered in Refs.~\cite{Fujikura2023,Bringmann2023}. These papers support the plausibility of MeV-to-sub-GeV dark-sector FOPTs as PTA-scale sources, but they are not the exact non-Abelian scalar-doublet model analyzed here.

\section{FOPT Results}
The FOPT agent used a semianalytic conformal high-temperature/polynomial backend with percolation and perturbativity checks. The scan ranges were
\[
0.25\le g_D\le 3.2,\qquad 10^{-5}\gev\le v_D\le 0.8\gev.
\]
Out of 934 scanned points, 662 were viable, 268 failed, and 4 were excluded by validation criteria. Representative high-ranking steering points are shown in Table~\ref{tab:fopt}.

\begin{center}
\captionof{table}{Representative FOPT steering points. Frequencies are order-of-magnitude steering estimates, not template-derived spectral peaks.}
\label{tab:fopt}
\begin{tabular}{lrrrrrr}
\toprule
Point & $g_D$ & $v_D$ [GeV] & $\alpha$ & $\beta/H$ & $T_n$ [GeV] & $f_{\rm est}$ [Hz]\\
\midrule
point\_0594 & 0.900 & $4.20\times10^{-2}$ & $2.10\times10^{10}$ & 21.21 & $1.56\times10^{-5}$ & $1.00\times10^{-8}$\\
point\_0769 & 1.400 & $4.11\times10^{-3}$ & $1.04\times10^{1}$ & 148.63 & $3.30\times10^{-4}$ & $9.92\times10^{-9}$\\
point\_0681 & 1.150 & $1.20\times10^{-2}$ & $9.30\times10^{2}$ & 62.18 & $2.86\times10^{-4}$ & $9.90\times10^{-9}$\\
point\_0441 & 1.823 & $1.04\times10^{-3}$ & $8.65\times10^{-1}$ & 381.21 & $2.00\times10^{-4}$ & $1.00\times10^{-8}$\\
point\_0715 & 1.250 & $7.03\times10^{-3}$ & $7.25\times10^{1}$ & 97.70 & $3.21\times10^{-4}$ & $9.81\times10^{-9}$\\
\bottomrule
\end{tabular}
\end{center}

The FOPT result is encouraging as a PTA-band steering calculation, but it is not a full finite-temperature effective-potential and bounce validation. The inherited caveats are the high-temperature approximation, Gaussian percolation approximation, lack of quantified gauge/daisy/RG uncertainty, and no independent wall-dynamics calculation.

\section{PTA Results}
The PTA stage used a benchmark scan plus a PTArcade campaign. For the benchmark comparison, the selected implemented spectral template was DBF. The scan considered 2611 points, of which 2546 were viable, and the best direct benchmark scored $13/14$ against the PTA violin windows. The best benchmark point was
\[
g_D=0.95,\qquad v_D=0.360438\gev,\qquad
\alpha=3.73\times10^8,\qquad \beta/H=23.69,\qquad T_\star=0.03237\gev.
\]
Its peak grid point was near $10^{-8}$ Hz with $h^2\Omega_{\rm GW}\simeq2.83\times10^{-9}$.

The PTArcade run status was \texttt{ok}. The campaign used the NANOGrav 15-year \texttt{ceffyl} mode with priors
\[
g_D\sim U(0.45,1.55),\qquad \log_{10}(v_D/\gev)\sim U(-3.2,0.7),
\]
and fixed DBF template. The accepted chain contains $N_{\rm samples}=100000$. The posterior mean values are
\[
g_D=0.8836\pm0.0353,\qquad \log_{10}(v_D/\gev)=-0.0539\pm0.1716,
\]
with $68\%$ credible intervals $g_D\in[0.8528,0.9159]$ and $\log_{10}(v_D/\gev)\in[-0.2217,0.1005]$. The MAP point is
\[
g_D=0.863175,\quad \log_{10}(v_D/\gev)=0.044548,\quad v_D=1.10802\gev,
\]
\[
T_{\rm reh}=0.09782\gev,\quad \alpha=2.14\times10^{14},\quad \beta/H=13.98,\quad f_{\rm peak,est}=1.75\times10^{-7}\,{\rm Hz}.
\]
The MAP $f_{\rm peak,est}$ is an approximate steering diagnostic and lies above the nominal PTA steering band; the fitted object is the DBF spectrum, not this scalar diagnostic alone.

\begin{figure}[htbp]
\centering
\includegraphics[width=0.82\linewidth]{pta_benchmark_spectrum.pdf}
\caption{Benchmark DBF spectrum compared with PTA violin windows. The best benchmark score is $13/14$.}
\label{fig:pta_benchmark}
\end{figure}

\begin{figure}[htbp]
\centering
\includegraphics[width=0.82\linewidth]{pta_ptarcade_spectrum.pdf}
\caption{PTArcade MAP DBF spectrum for the accepted widened-prior run.}
\label{fig:pta_ptarcade}
\end{figure}

\begin{figure}[htbp]
\centering
\includegraphics[width=0.82\linewidth]{pta_posteriors.pdf}
\caption{PTArcade posterior distributions for the fitted parameter basis.}
\label{fig:pta_posteriors}
\end{figure}

The main PTA caveat is template validity. The DBF template was selected as the closest implemented template for very large $\alpha$, while the Higgsless template was not selected because many local points exceed its stated $\alpha\lesssim0.5$ validation hint. DBF amplitudes use backend defaults, including efficiency assumptions, while wall velocity and microscopic efficiency are not computed independently.

\section{Constraints}
The constraints agent considered 12 constraint classes: 1 directly testable theory check, 7 requiring new backend calculations, and 4 not applicable to the exact zero-portal model. The directly testable perturbativity check passed for the posterior region, with $g_D$ well below $\sqrt{4\pi}$. This is a basic consistency check, not a complete theory-validity proof.

\begin{center}
\captionof{table}{Constraint classification. ``No exclusion'' means no exclusion with the observables currently available.}
\label{tab:constraints}
\begin{tabular}{p{0.16\linewidth}p{0.24\linewidth}p{0.20\linewidth}p{0.28\linewidth}}
\toprule
Class & Examples & Applicability & Missing ingredient \\
\midrule
BBN/CMB thermal dark sectors & Refs.~\cite{Sabti2019,Sabti2021} & Requires new backend & $T_D/T_{\rm SM}$, abundance, entropy transfer, likelihood map. \\
$\Delta N_{\rm eff}$ and dark radiation & Planck~\cite{Planck2018} & Requires new backend & Decoupling history and radiation budget. \\
PTA-scale energy release & Ref.~\cite{Bringmann2023} & Requires new backend & Visible/dark reheating partition and late energy injection. \\
Self-interacting DM & Ref.~\cite{TulinYu2018}; Bullet Cluster~\cite{Markevitch2004} & Requires new backend & $\sigma_T/m(v)$ and dark-vector DM fraction. \\
Portal laboratory searches & Refs.~\cite{Arcadi2019,Alexander2016,Battaglieri2017} & Not applicable & Zero portal in the minimal model; relevant again for portal extensions. \\
Perturbativity & Upstream handoffs & Direct & Passed basic $g_D<\sqrt{4\pi}$ check; higher-loop validity not tested. \\
\bottomrule
\end{tabular}
\end{center}

The most important missing backend calculations are a thermal-history module, BBN/CMB observable mapping, vector relic-density calculation, non-Abelian vector self-interaction cross sections, portal widths for non-minimal variants, and higher-order/RG validity diagnostics.

\section{Assumptions and Limitations}
The prior audit identified 18 assumptions, limitations, priors, and caveats. The highest-severity items are:
\begin{itemize}
\item No thermal history, dark-sector temperature ratio, decoupling history, reheating partition, relic abundance, or entropy-transfer observable was computed.
\item Sub-GeV strongly supercooled FOPTs can affect BBN/CMB through energy release or dark radiation, but $\Delta N_{\rm eff}$ and visible/dark reheating were not mapped.
\item The minimal model fixes the Higgs portal to zero, which removes direct portal searches but leaves the origin and thermal population of the dark sector unspecified.
\item No dark-vector relic abundance or fractional dark-matter abundance was computed, so dark-matter viability and self-interaction constraints cannot be applied.
\item The PTA analysis uses DBF spectra at extremely large $\alpha$ without an independent calculation of wall velocity or microscopic efficiency.
\item The FOPT calculation is semianalytic and does not quantify gauge, daisy, RG-scale, or full-bounce uncertainties.
\end{itemize}
Consequently, the PTA fit should be interpreted as a conditional spectrum-level match within the implemented backend, not as a complete cosmological viability proof.

\section{Conclusions}
Robust claims are limited to the structure of the local candidate and the executed pipeline outputs. The minimal conformal non-Abelian candidate is a dark \SUtwoD{} gauge theory with one complex scalar doublet, classically scale-invariant tree potential, and two local free parameters $(g_D,v_D)$. The exact core model has same-model precedent in Ref.~\cite{Borah2021}. The local FOPT scan found many viable points, and the PTA workflow found a strong benchmark near-match and an \texttt{ok} PTArcade posterior run.

Conditional claims are the PTA interpretation and parameter preference. They depend on the DBF template, backend approximations, chosen priors, and missing wall-dynamics inputs. The reported MAP point is a valid output of the local PTArcade campaign, but its extremely large $\alpha$ and approximate peak diagnostic require caution.

Blocked claims are cosmological viability, dark-matter viability, and self-interaction safety. These require new thermal-history, BBN/CMB, $\Delta N_{\rm eff}$, relic-density, and self-interaction backends.

Uncontrolled claims would include saying that the model is novel, proven viable, or excluded-free in a full cosmological sense. The current constraints result is better phrased as: no exclusion was derived from the available observables, and most relevant observables remain missing.

\begin{thebibliography}{99}
\bibitem{Borah2021}
D. Borah, A. Dasgupta and S. K. Kang,
``A first order dark $SU(2)_D$ phase transition with vector dark matter in the light of NANOGrav 12.5 yr data,''
JCAP 12 (2021) 039,
arXiv:2109.11558.

\bibitem{Balan2025}
S. Balan, T. Bringmann, F. Kahlhoefer, J. Matuszak and C. Tasillo,
``Sub-GeV dark matter and nano-Hertz gravitational waves from a classically conformal dark sector,''
JCAP 08 (2025) 062,
arXiv:2502.19478.

\bibitem{Fujikura2023}
K. Fujikura, S. Girmohanta, Y. Nakai and M. Suzuki,
``NANOGrav Signal from a Dark Conformal Phase Transition,''
arXiv:2306.17086.

\bibitem{Goncalves2025}
J. Goncalves, D. Marfatia, A. P. Morais and R. Pasechnik,
``Supercooled phase transitions in conformal dark sectors explain NANOGrav data,''
Phys. Lett. B 869 (2025) 139829,
arXiv:2501.11619.

\bibitem{Bringmann2023}
T. Bringmann, P. F. Depta, T. Konstandin, K. Schmidt-Hoberg and C. Tasillo,
``Does NANOGrav observe a dark sector phase transition?''
arXiv:2306.09411.

\bibitem{Sabti2019}
N. Sabti, J. Alvey, M. Escudero, M. Fairbairn and D. Blas,
``Refined Bounds on MeV-scale Thermal Dark Sectors from BBN and the CMB,''
arXiv:1910.01649.

\bibitem{Sabti2021}
N. Sabti, J. Alvey, M. Escudero, M. Fairbairn and D. Blas,
``Addendum: Refined bounds on MeV-scale thermal dark sectors from BBN and the CMB,''
arXiv:2107.11232.

\bibitem{Planck2018}
Planck Collaboration,
``Planck 2018 results. VI. Cosmological parameters,''
arXiv:1807.06209.

\bibitem{TulinYu2018}
S. Tulin and H.-B. Yu,
``Dark Matter Self-interactions and Small Scale Structure,''
Phys. Rept. 730 (2018) 1,
arXiv:1705.02358.

\bibitem{Markevitch2004}
M. Markevitch et al.,
``Direct constraints on the dark matter self-interaction cross-section from the merging galaxy cluster 1E0657-56,''
Astrophys. J. 606 (2004) 819,
arXiv:astro-ph/0309303.

\bibitem{Arcadi2019}
G. Arcadi, A. Djouadi and M. Raidal,
``Dark Matter through the Higgs portal,''
Phys. Rept. 842 (2020) 1,
arXiv:1903.03616.

\bibitem{Alexander2016}
J. Alexander et al.,
``Dark Sectors 2016 Workshop: Community Report,''
arXiv:1608.08632.

\bibitem{Battaglieri2017}
M. Battaglieri et al.,
``US Cosmic Visions: New Ideas in Dark Matter 2017: Community Report,''
arXiv:1707.04591.
\end{thebibliography}

\end{document}
